\documentclass[letterpaper,10pt]{article}
\usepackage[utf8]{inputenc}

%% PACKAGES

\nofiles
\usepackage{titlesec}
\usepackage[usenames,dvipsnames]{color}
\usepackage[empty]{fullpage}
\usepackage{geometry}
\usepackage{hyperref}
\usepackage{multirow}
\usepackage[T1]{fontenc}
\usepackage{color}
\usepackage{array}
\usepackage{fontawesome}
\usepackage{hyperref}
\usepackage{colortbl}
\usepackage{enumitem}

%% MARGINS %%

\usepackage{graphicx}
\newcommand\sbullet[1][.5]{\mathbin{\vcenter{\hbox{\scalebox{#1}{$\bullet$}}}}}

\geometry{
    letterpaper,
    left=12mm,
    top=12mm, % Increased slightly for better header spacing
    right=12mm,
    bottom=12mm,
}

\raggedright

%% CUSTOM COMMANDS %%

% Color on lines and name
\definecolor{burgundy}{rgb}{0,0,0}
\definecolor{darkblue}{rgb}{0,0,0.5}
\definecolor{light-grey}{rgb}{0.4,0.4,0.4}

\hypersetup{colorlinks=true,urlcolor=darkblue}

% Command for displaying organizations (schools & employers)
\newcommand{\organization}[4]{
    \vspace{-1.0pt}
    \textbf{#1} \hfill{\emph{#2}} \\
    \emph{#3} \hfill{#4} \\
    \vspace{1.5pt} % Reduced from 3pt
}

% Command for the beginning of bullets
\newcommand{\bulletsBegin}{
    \vspace{1pt}
    \begin{minipage}{17.6cm} 
    \begin{itemize}[leftmargin=0.6cm]
    \setlength\itemsep{1pt} % Increased spacing slightly for readability
}

% Command for the end of bullets
\newcommand{\bulletsEnd}{
    \end{itemize}\vspace{1pt}
    \end{minipage}
}

% Command for email
\newcommand{\myEmail}[1]{
    \href{mailto:#1}{\faEnvelope \ #1}
    \ |
}

% Command for github profile
\newcommand{\myGitHub}[1]{
    \href{https://github.com/#1/}{\faGithub \ #1}
    \ 
}

% Command for displaying name and Contact info
\newcommand{\myName}[5]{
    \begin{center}
        {\huge{\color{burgundy}\scshape{#1}}} \\
        \vspace{5pt}
        {\large \scshape #5} \\ % This is the new Title/Location line
        \vspace{5pt}
        \myEmail{#4}
        \href{https://peterzeng.github.io/}{\faLink \ peterzeng.github.io} \ | \ 
        \myGitHub{#3}
    \end{center}
    \vspace{-4pt}
}

% Title formatting for each section
\titleformat{\section}{
    \vspace{-10pt}\color{burgundy}\scshape\raggedright\large
    }{}{0em}{}[\titlerule]

%% DOCUMENT 

\begin{document}

    %% HEADING --------------------
    \myName{Peter Zeng}{}{peterzeng}{peterateast@gmail.com}
    {PhD Candidate in Computer Science $\sbullet[.75]$ Stony Brook, NY} 

    %% EDUCATION --------------------
    \section{Education}
    \organization{Stony Brook University}{August 2023 - Expected May 2027}
        {Ph.D. in Computer Science}{Advisor: \href{https://owenrambow.com/}{Owen Rambow}} 

    \organization{Stony Brook University}{August 2017 - May 2022}
        {BS/MS in Computer Science}{}

    %% RESEARCH INTERESTS --------------------
    \section{Research Interests}
    \vspace{-2pt}
    \begin{itemize}[leftmargin=0.6cm]
        % \setlength\itemsep{2pt}
        \item \textbf{Machine Theory of Mind \& Cognitively Informed AI:} Assessing the ability of Large Language Models (LLMs) and Vision-Language Models (LVLMs) to model belief states, track recursive mental states, and perform human-like cognitive reasoning.
        \item \textbf{Multimodal Grounding \& Human-AI Communication:} Investigating how LVLMs ground language to vision and coordinate with human partners, with a focus on divergences between human pragmatic inference and model behavior in interactive settings.
        \item \textbf{AI Interpretability \& Trustworthiness:} Developing hybrid architectures that combine the performance of deep neural networks with the transparency of feature-based models (e.g., residualized similarity).
    \end{itemize}
    % \vspace{8pt}

    %% RESEARCH EXPERIENCE (PUBLICATIONS) --------------------
    \section{Selected Publications \& Manuscripts}

    \organization{\href{https://arxiv.org/pdf/2606.17372}{Implicit vs. Explicit Prompting Strategies for LVLMs in Referential Communication}}{Spring 2026}
    {\textbf{*Peter Zeng}, Amie J. Paige, Weiling Li, Susan E. Brennan, Owen Rambow, Cameron R. Jones. \textit{Under Review}.}{}
    \bulletsBegin
        \item \textbf{Problem:} Two prior studies reached contradictory conclusions about whether LVLMs can coordinate efficient referring expressions, leaving it unclear whether task design or prompting approach was responsible.
        \item \textbf{Method:} Controlled for task differences between the conflicting studies while directly comparing explicit versus implicit prompting strategies across models, isolating prompting as the key variable.
        \item \textbf{Outcome:} Showed that LVLMs succeed at communicative efficiency when explicitly instructed to do so, but fail to infer this need from implicit prompts — unlike humans, who derive communicative intent from context alone — revealing a fundamental gap in pragmatic inference.
    \bulletsEnd

    \organization{\href{https://arxiv.org/abs/2601.19792v2}{LVLMs and Humans Ground Differently in Referential Communication}}{Fall 2025}
    {\textbf{*Peter Zeng}, Weiling Li, Amie J. Paige, Zhengxiang Wang, Panagiotis Kaliosis, \\ Dimitris Samaras, Gregory Zelinsky, Susan E. Brennan, and Owen Rambow. \textit{ACL 2026 (Main)}.}{}
    \bulletsBegin
        \item \textbf{Problem:} LVLMs appear competent in multimodal reasoning, but it's unclear whether they ground the same way humans do during interactive communication.
        \item \textbf{Method:} Designed and implemented a referential communication task comparing human–human and human–LVLM interactions. Analyzed linguistic strategies, visual attention patterns, and error profiles across conditions to isolate model–human divergences.
        \item \textbf{Outcome:} Demonstrated that LVLMs systematically rely on surface‑level visual heuristics and fail to adapt grounding strategies to partner feedback, unlike humans who use pragmatic inference and incremental clarification. Findings reveal a fundamental mismatch between human and model grounding processes and establish a new evaluation paradigm for interactive multimodal communication.
    \bulletsEnd
    
    % PAPER 1
    \organization{\href{https://aclanthology.org/2025.findings-emnlp.856/}{Residualized Similarity for Interpretability in Authorship Verification}}{Spring 2025}
    {\textbf{*Peter Zeng}, Pegah Alipoormolabashi, Jihu Mun, Gourab Dey, Nikita Soni, \\Niranjan Balasubramanian, Owen Rambow, H. Schwartz. \textit{EMNLP 2025 (Findings)}.}{}
    \bulletsBegin
        \item \textbf{Problem:} Standard cosine-similarity models in authorship attribution often sacrifice interpretability for performance.
        \item \textbf{Method:} Developed a novel neural framework that predicts the "residual" (error) of an interpretable baseline, allowing the model to capture complex non-linearities while retaining feature-level explainability.
        \item \textbf{Outcome:} Improved accuracy by $\sim$20-30\% over baseline interpretable systems, establishing a new SOTA for faithful explainability in authorship tasks.
    \bulletsEnd

    % PAPER 2
    \organization{\href{https://arxiv.org/abs/2406.12131}{Gram2Vec: An Interpretable Document Vectorizer}}{}
        {\textbf{*Peter Zeng}, Hannah Stortz, Eric Sclafani, Alina Shabaeva,\\ Maria Elizabeth Garza, Daniel Greeson, Owen Rambow. \textit{Under Review}.}{}
        \bulletsBegin
            \item \textbf{Problem:} Neural text embeddings are often "black boxes," making it difficult to understand the linguistic basis for authorship attribution.
            \item \textbf{Method:} Created a grammatical style embedding algorithm that extracts normalized relative frequencies of morphosyntactic features into a higher-dimensional space.
            \item \textbf{Outcome:} Demonstrated that explicit grammatical feature extraction can be successfully applied to two domains: authorship verification and AI detection.
        \bulletsEnd
    
    % PAPER 3
    \organization{\href{https://arxiv.org/pdf/2502.06922}{Synthetic Audio Helps for Cognitive State Tasks}}{Fall 2024}
    {Adil Soubki, John Murzaku, \textbf{*Peter Zeng}, Owen Rambow. \textit{NAACL 2025} (Findings).}{}
    \bulletsBegin
        \item \textbf{Problem:} Gold-standard audio data for cognitive tasks is scarce, yet text-only models miss prosodic cues.
        \item \textbf{Method:} Introduced Synthetic Audio Data fine-tuning (SAD), a multimodal training approach using text and zero-shot synthetic audio from off-the-shelf TTS systems.
        \item \textbf{Outcome:} Proved that synthetic audio acts as an effective regularizer, improving cognitive state modeling performance even when real audio is unavailable.
    \bulletsEnd
    
    % PAPER 4
    \organization{\href{https://arxiv.org/pdf/2403.02451}{Views Are My Own, But Also Yours: Benchmarking Theory of Mind}}{Fall 2023}
        {Adil Soubki, John Murzaku, Arash Yousefi Jordehi, \textbf{*Peter Zeng},\\ Magdalena Markowska, Seyed Abolghasem Mirroshandel, Owen Rambow. \textit{ACL 2024} (Findings).}{}
        \bulletsBegin
            \item \textbf{Problem:} Existing Theory of Mind benchmarks rely on synthetic scenarios that do not reflect the complexity of natural dialogue.
            \item \textbf{Method:} Constructed \textit{Common-ToM}, a benchmark derived from naturally occurring dialogs, to test LLM capabilities zero-shot.
            \item \textbf{Outcome:} Demonstrated that while LLMs retrieve facts well, they struggle to track recursive belief states without explicit belief representations.
        \bulletsEnd
        
    % PAPER 5
    \organization{\href{https://aclanthology.org/2022.coling-1.66}{Re-Examining FactBank: Predicting the Author's Presentation of Factuality}}{Spring 2022}
        {John Murzaku, \textbf{*Peter Zeng}, Magdalena Markowska, Owen Rambow. \textit{COLING 2022}.}{}
        \bulletsBegin
            \item \textbf{Problem:} Prior research on the canonical FactBank corpus relied on data representation errors, leading to invalid conclusions.
            \item \textbf{Method:} Systematically identified and corrected representation errors in the corpus to create a reliable ground truth.
            \item \textbf{Outcome:} Achieved State-of-the-Art results for predicting authorial factuality across four distinct corpora using the corrected data.
        \bulletsEnd

    %% PROFESSIONAL EXPERIENCE --------------------
    \section{Professional Experience}

    \organization{Applied Scientist Intern}{June 2026 - September 2026}
        {\href{https://business.adobe.com/products/experience-platform/adobe-experience-platform.html}{Adobe (AEP Agent Orchestrator)}}{San Jose, CA}
        \bulletsBegin
            \item Designing and implementing memory management systems for the Adobe Experience Platform (AEP) Agent Orchestrator, enabling agents to persist, retrieve, and reason over long-horizon context.
            \item Researching memory architectures (episodic, semantic, and working memory) to improve agent coherence and task continuity across multi-step orchestration workflows.
            \item Collaborating with cross-functional teams to evaluate memory retrieval strategies and their downstream impact on agent response quality and latency.
        \bulletsEnd

    \organization{Research Assistant}{June 2025 - Present}
        {\href{https://www.stonybrook.edu/commcms/bias-nrt/}{BIAS-NRT (NSF Fellow)}}{Stony Brook, NY}
        % Updated description based on Susan's notes
        \emph{Detecting and Addressing Bias in Data, Humans, \& Institutions}
        \bulletsBegin
            \item Conducting task analyses with stakeholders at an Innocence Organization to define user needs and iteratively design optimal prompts for CoCounsel (a secure LLM) to assist in the intake process.
            \item Developing multimodal evaluation frameworks to assess how Vision-Language Models (VLMs) describe and match objects compared to human baselines.
            \item Investigating bias propagation in large-scale datasets used for training generative models.
        \bulletsEnd

    \organization{Research Assistant}{August 2023 - May 2025}
        {\href{https://www.iarpa.gov/research-programs/hiatus}{IARPA HIATUS (Stony Brook University)}}{Stony Brook, NY}
        \bulletsBegin
            \item Lead researcher on the "Residualized Similarity" project, focusing on attribution in low-resource and obfuscated text domains.
            \item Managed data pipelines for large-scale authorship corpora, ensuring rigorous train/test separation to prevent data leakage in open-world attribution tasks.
        \bulletsEnd
        
    \organization{Software Engineer}{June 2022 - August 2023}
        {\href{https://aws.amazon.com/ecs/}{Amazon Web Services (Elastic Container Service)}}{New York, NY}
        \bulletsBegin
            \item \textbf{High-Scale Infrastructure:} Maintained core control-plane infrastructure for Amazon ECS, supporting 2 million daily customer deployments across global regions.
            \item \textbf{Reliability Engineering:} Designed and implemented alarm-based rollback mechanisms, requiring the ingestion and real-time processing of massive-scale telemetry data to ensure service availability.
            \item \textbf{Production Excellence:} Participated in on-call rotations for a Tier-1 AWS service, resolving critical production incidents affecting enterprise customers.
        \bulletsEnd
        
    %% TECHNICAL SKILLS (NEW) --------------------
    \section{Technical Skills}
    \vspace{-2pt}
    \begin{tabular}{@{}p{5cm}l}
        \textbf{Languages} & Python, Java \\
        \textbf{Deep Learning Libraries} & PyTorch, HuggingFace (Transformers, PEFT), W\&B \\
        \textbf{AI Techniques} & Fine-Tuning, Prompt Engineering, RAG \\
        \textbf{Infrastructure} & Git, Linux/Unix \\
        \textbf{Methodologies} & Causal Inference, Human-Subjects Evaluation, Survey Design \\
    \end{tabular}
    \vspace{8pt}

    %% ACADEMIC SERVICE & TEACHING (UPDATED) --------------------
    \section{Academic Service \& Teaching}
    \vspace{-2pt}
    \begin{itemize}[leftmargin=0.6cm]
        \setlength\itemsep{2pt}
        \item \textbf{Reviewer:} ACL Rolling Review (ARR), NeurIPS 2026 Evaluations and Datasets Track
        \item \textbf{Teaching Assistant (Stony Brook):} CSE 385 Analysis of Algorithms, CSE 160/260 OOP \& Data Structures
        \item \textbf{Academic Tutor (ASTC):} CSE 101 (Intro to Programming) \& PHY 131/132 (Physics Mechanics) at the \emph{Academic Success and Tutoring Center}, Stony Brook, NY
        \item \textbf{Instructor:} Taught K-12 introductory programming (Scratch, Python, Java) at \emph{theCoderSchool}, Buffalo, NY
    \end{itemize}

\end{document}